\chapter{Homomorphisms}

\section{Homomorphisms and Isomorphisms}

\begin{definition}
  A map \(\f:G\to G'\) is called a \emph{group homomorphism} if \[
    \f(ab) = \f(a)\f(b) \text{ for all } a,b\in G.
  \]
  If \(\f\) is bijective, then it is called a \emph{group isomorphism}. In this case, we write \(G\cong G'\). Further, if \(G' = G\), then \(\f\) is called a \emph{group automorphism}.
\end{definition}

\begin{example}
  The determinant \(\det:GL_n(\R)\to \R^\times\) is a homomorphism.
\end{example}

\begin{example}
  The exponential function \(\exp:(\R,+)\to (\R_{>0}, \cdot)\) is an isomorphism with inverse \(\log\).
\end{example}

\begin{example}
  For \(n\neq 0\), we have \(n\Z\leq\Z\) and the map \(\f:n\Z\to\Z\) defined by \[
    \f(nk) = k
  \] is an isomorphism.
\end{example}

\begin{example}
  For \(n > 0\), the map \(\f:\Z\to\Z_n\) defined by \[
    \f(k) = k\mod n
  \] is a homomorphism \(\ker\f = n\Z\).
\end{example}

\begin{example}
  The map \[
    SO_2(\R)\to U(1), \quad
    \begin{pmatrix}
      \cos\theta & -\sin\theta \\
      \sin\theta & \cos\theta
    \end{pmatrix}\mapsto e^{i\theta}
  \] is an isomorphism.
\end{example}

\begin{example}
  The finite groups \(\Z_4\) and \(\Z_2\times Z_2\) are \emph{not} isomorphic although they have the same order.
\end{example}
\begin{proof}
  Recall that \[
    \Z_4 = \{0,1,2,3\}, \quad
    \Z_2\times\Z_2 = \{(0,0),(0,1),(1,0),(1,1)\}.
  \]
  \begin{itemize}
    \item In \(Z_4\), \(|1| = 4\).
    \item In \(\Z_2\times\Z_2\), \(|(1,0)| = |(0,1)| = |(1,1)| = 2\).
  \end{itemize}
  Suppose \(\f:\Z_4\to\Z_2\times\Z_2\) is a homomorphism, then \[
    \f(2) = 2\f(1) = (0,0),
  \] which forces \(\f\) to be non-injective.
\end{proof}

\section{Cyclic Subgroups}

\section{Cyclic Groups}
\begin{definition}
  Let \(G\) be a group and \(a\in G\). Then \[
    \<a\> := \{a^n \mid n\in\Z\}
  \] is a subgroup of \(G\) called the \emph{cyclic subgroup generated by \(a\)}. If \(G = \<a\>\), then \(G\) is called a \emph{cyclic group}.
\end{definition}

\begin{proposition}
  Every cyclic group is abelian.
\end{proposition}
\begin{proof}
  Suppose \(G = \<a\>\), let \(x = a^m\) and \(y = a^n\), then \[
    xy = a^ma^n = a^{m+n} = a^{n+m} = a^na^m = yx.\qedhere
  \]
\end{proof}

\begin{theorem}[Classification of cyclic groups]
  Any infinite cyclic group is isomorphic to \((\Z,+)\) and any finite cyclic group of order \(n\) is isomorphic to \((\Z_n,+)\).
\end{theorem}

\begin{theorem}[Classification of subgroups of a finite cyclic group]
  Let \(G = \<a\>\) be a cyclic group of finite order \(n\). Let \(a^s\in G\), then \[
    |a^s| = \frac{n}{\gcd(s,n)}.
  \] Moreover, \(\<a^s\> = \<a^t\>\) if and only if \(\gcd(s,n) = \gcd(t,n)\).
\end{theorem}

\section{Generating Sets}